% ---- pg-457 （书页 445） ----

且令向量 $\mathbf{X}^{(2)}=[X_1^{(2)},X_2^{(2)}]'$ 的变量代表刚性的静止测度, 对这些数据作典型相关分析.

\noindent 10.16\quad 安德鲁斯和赫茨伯格(Andrews and Herzberg, 见文献[1])给出由比较无糖尿病和有糖尿病的病人得到的数据. 三个主要变量

$X_1^{(1)}=$ 葡萄糖允许量

$X_2^{(1)}=$ 胰岛素对口部葡萄糖的反应

$X_3^{(1)}=$ 胰岛素抵抗力

和两个次要变量

$X_1^{(2)}=$ 相对重量

$X_2^{(2)}=$ 禁食的血浆葡萄糖

都被测量了. 从 $n=46$ 个无糖尿病病人的数据得到协方差矩阵
\[
\mathbf{S}=\left[\begin{array}{c:c}\mathbf{S}_{11}&\mathbf{S}_{12}\\ \hdashline \mathbf{S}_{21}&\mathbf{S}_{22}\end{array}\right]=\left[\begin{array}{ccc:cc}
1\,106.000&396.700&108.400&0.787&26.230\\
396.700&2\,382.000&1\,143.000&-0.214&-23.960\\
108.400&1\,143.000&2\,136.000&2.189&-20.840\\ \hdashline
0.787&-0.214&2.189&0.016&0.216\\
26.230&-23.960&-20.840&0.216&70.560
\end{array}\right]
\]
求样本典型变量及其相关系数, 解释这些量. 第一典型变量分别恰当地综合了它们各自变量组吗? 试说明. 检验典型相关系数在 $\alpha=0.05$ 时的显著性.

\noindent 10.17\quad 关于人对吸烟的渴望以及心理状态和身体状态的数据, 从 $n=110$ 个试验对象收集而来. 数据是对 12 个问题(变量)的回答, 编码为 1 到 5. 与对吸烟的渴望有关的四个标准化测度量定义如下

$\mathbf{z}_1^{(1)}=$ 吸烟 1(第一措词)

$\mathbf{z}_2^{(1)}=$ 吸烟 2(第二措词)

$\mathbf{z}_3^{(1)}=$ 吸烟 3(第三措词)

$\mathbf{z}_4^{(1)}=$ 吸烟 4(第四措词)

与心理和身体状态有关的八个标准测量值定义如下

$\mathbf{z}_1^{(2)}=$ 集中力

$\mathbf{z}_2^{(2)}=$ 烦恼

$\mathbf{z}_3^{(2)}=$ 睡眠

$\mathbf{z}_4^{(2)}=$ 紧张

$\mathbf{z}_5^{(2)}=$ 警惕

$\mathbf{z}_6^{(2)}=$ 急躁

$\mathbf{z}_7^{(2)}=$ 疲劳

$\mathbf{z}_8^{(2)}=$ 满意

由这些数据构造的相关系数矩阵为
\[
\mathbf{R}=\left[\begin{array}{c:c}\mathbf{R}_{11}&\mathbf{R}_{12}\\ \hdashline \mathbf{R}_{21}&\mathbf{R}_{22}\end{array}\right]
\]
其中
